% ==============================================================================
% inner_cover.tex — Academic Dark-Mode Vector Art Inner Title Page
% Standards: academic-book-cover-architect (TikZ Bezier Curves & Pure Vector 100%)
% ==============================================================================
\begin{titlepage}
\thispagestyle{empty}
\pagecolor{black!95}

\begingroup
\setlength{\XeTeXlinebreakskip}{0pt}

\begin{tikzpicture}[remember picture, overlay]
    % 1. Subtle Cybernetic Background Grid
    \draw[cyan!20!black, line width=0.4pt, step=0.4in] 
         ([shift={(0.2in,0.2in)}]current page.south west) grid ([shift={(-0.2in,-0.2in)}]current page.north east);

    % 2. Central Vector Holographic Rings (Pure TikZ)
    \draw[rbruEmerald!60!white, line width=1.5pt, opacity=0.4] 
         ([yshift=0.1in]current page.center) circle (1.8in);
    \draw[rbruGold!80!white, line width=0.8pt, opacity=0.5, dashed] 
         ([yshift=0.1in]current page.center) circle (2.1in);
    \draw[cyan!70!white, line width=0.6pt, opacity=0.3] 
         ([yshift=0.1in]current page.center) circle (2.3in);

    % Hexagonal Vector Accent
    \draw[rbruGold, line width=1pt, opacity=0.7]
         ([yshift=0.1in]current page.center) ++(0:1.8in)
         -- ([yshift=0.1in]current page.center) ++(60:1.8in)
         -- ([yshift=0.1in]current page.center) ++(120:1.8in)
         -- ([yshift=0.1in]current page.center) ++(180:1.8in)
         -- ([yshift=0.1in]current page.center) ++(240:1.8in)
         -- ([yshift=0.1in]current page.center) ++(300:1.8in) -- cycle;

    \foreach \i in {0,60,120,180,240,300} {
        \draw[rbruEmerald!80!white, fill=black, line width=1pt] 
             ([yshift=0.1in]current page.center) ++(\i:1.8in) circle (4pt);
    }

    % Glowing Central Core Light
    \fill[rbruEmerald!40!white, opacity=0.15] ([yshift=0.1in]current page.center) circle (1.2in);

    % 3. Header & Document Type
    \node[anchor=north, align=center, text width=5.2in] 
          at ([yshift=-0.7in]current page.north) {
        {\fontsize{9}{11}\selectfont\color{rbruEmerald!70!white}%
        \textbf{SERIES OF AI SKILL \textbullet\ VOLUME 05: CLOUD DATA \& ML}}\\[3pt]
        {\fontsize{8}{10}\selectfont\color{gray!40}%
        ศูนย์ความเป็นเลิศด้านสิ่งแวดล้อมและพลังงาน (LEQs) มหาวิทยาลัยราชภัฏรำไพพรรณี}
    };

    % 4. Main Title & Subtitle (Vector Typography)
    \node[anchor=north, align=center, text width=5.6in] 
          at ([yshift=-1.3in]current page.north) {
        {\fontsize{17}{22}\selectfont\bfseries\color{white}%
        ปัญญาประดิษฐ์กับวิศวกรรมข้อมูลคลาวด์\\
BigQuery และการเรียนรู้ของเครื่อง}\\[8pt]
        {\fontsize{8.5}{11}\selectfont\color{gray!30}%
        AI FOR CLOUD DATA ENGINEERING, BIGQUERY \& MACHINE LEARNING}
    };

    % 5. Author & Academic Information
    \node[anchor=south, align=center, text width=5.2in] 
          at ([yshift=0.7in]current page.south) {
        {\fontsize{13}{16}\selectfont\bfseries\color{white}%
        ผู้ช่วยศาสตราจารย์ ดร.ชีวะ ทัศนา}\\[4pt]
        {\fontsize{9}{12}\selectfont\color{gray!30}%
        สาขาวิชาฟิสิกส์ คณะวิทยาศาสตร์และเทคโนโลยี}\\[3pt]
        {\fontsize{8.5}{11}\selectfont\color{gray!40}%
        ศูนย์ความเป็นเลิศด้านสิ่งแวดล้อมและพลังงาน (LEQs) มหาวิทยาลัยราชภัฏรำไพพรรณี}\\[8pt]
        {\color{gray!50}\rule{3.8in}{0.5pt}}\\[8pt]
        {\fontsize{8}{10}\selectfont\color{rbruGold!80!white}%
        โครงการจัดทำตำราวิชาการและชุดความรู้ปัญญาประดิษฐ์ \textbullet\ พ.ศ. 2569}
    };

\end{tikzpicture}

\endgroup
\end{titlepage}

\pagecolor{white}
\clearpage
\thispagestyle{empty}
\null\clearpage
